\DocumentMetadata{
  pdfversion=2.0,
  pdfstandard=ua-2,
  lang=en-US,
  tagging=on,
  tagging-setup={math/setup=mathml-SE,tabsorder=structure}
}
\documentclass[11pt,letterpaper]{article}
\usepackage[margin=0.68in,includefoot]{geometry}
\usepackage{newcomputermodern}
\usepackage{amsmath,amscd,mathtools}
\usepackage{tikz,adjustbox}
\providecommand{\square}{\mdlgwhtsquare}
\usepackage[hidelinks]{hyperref}
\hypersetup{
  pdftitle={Analysis PhD Exam, May 2008},
  pdfauthor={University of Florida Department of Mathematics},
  pdfsubject={Graduate mathematics examination},
  pdfdisplaydoctitle=true,
  bookmarksopen=true
}
\setcounter{secnumdepth}{0}
\setlength{\parindent}{0pt}
\setlength{\parskip}{0.28em}
\setlength{\emergencystretch}{3em}
\setlength{\leftmargini}{2.15em}
\setlength{\leftmarginii}{2.65em}
\setlength{\leftmarginiii}{3em}
\pagestyle{empty}
\raggedbottom
\allowdisplaybreaks
\NewTaggingSocketPlug{block/list/label}{examproblem}{%
  \tagstructbegin{tag=Lbl}\tagstructend%
  \tagstructbegin{tag=\LBody}%
  \tagstructbegin{tag=\ProblemHeadingTag,title-o={\ProblemHeadingText}}%
  \tagmcbegin{tag=\ProblemHeadingTag,actualtext={\ProblemHeadingText}}%
  #2%
  \tagmcend%
  \tagstructend%
}
\newenvironment{examproblems}[2]{%
  \def\ProblemHeadingTag{#2}%
  \AssignTaggingSocketPlug{block/list/label}{examproblem}%
  \begin{enumerate}%
  \setcounter{enumi}{#1}%
  \renewcommand{\labelenumi}{\textbf{\theenumi.}}%
  \setlength{\topsep}{0.35em}%
  \setlength{\partopsep}{0pt}%
  \setlength{\itemsep}{0.48em}%
  \setlength{\parsep}{0.18em}%
}{\end{enumerate}}
\newenvironment{examsubparts}{%
  \AssignTaggingSocketPlug{block/list/label}{default}%
  \begin{enumerate}%
  \setlength{\topsep}{0.18em}%
  \setlength{\partopsep}{0pt}%
  \setlength{\itemsep}{0.16em}%
  \setlength{\parsep}{0.1em}%
}{\end{enumerate}}
\newenvironment{examinstructions}{%
  \par\begingroup\itshape
}{%
  \par\endgroup\vspace{0.18em}
}
\makeatletter
\renewcommand\subsection{\@startsection{subsection}{2}{0pt}%
  {-1.25ex plus -.25ex minus -.1ex}%
  {0.55ex plus .1ex}%
  {\normalfont\large\bfseries}}
\renewcommand{\@maketitle}{%
  \newpage\null\vskip -1em%
  {\footnotesize\bfseries
    \href{https://www.ufl.edu/}{University of Florida}, \href{https://math.ufl.edu/}{Department of Mathematics}\par}%
  \vskip -0.5em%
  \tagstructbegin{tag=H1,title={Analysis PhD Exam, May 2008}}%
  \tagpdfparaOff%
  \tagmcbegin{tag=H1}%
    {\LARGE\bfseries\href{https://gma.math.ufl.edu/exams/analysis-phd/}{Analysis PhD Exam}, May 2008\par}%
  \tagmcend%
  \tagpdfparaOn%
  \tagstructend%
  \vskip 0.25em%
  \tagmcbegin{artifact}%
    \hrule height 1pt%
  \tagmcend%
  \vskip 1em%
}
\makeatother
\title{Analysis PhD Exam, May 2008}
\author{}
\date{}
\begin{document}
\maketitle
\thispagestyle{empty}
\begin{examinstructions}
Write solutions in a neat and logical fashion, giving complete reasons for all steps.
\end{examinstructions}
\begin{examproblems}{0}{H2}
\def\ProblemHeadingText{Problem 1.}
\item
Define the Hardy–Littlewood maximal function. State and prove the Hardy–Littlewood maximal theorem. (Begin by proving a suitable covering lemma.)
\def\ProblemHeadingText{Problem 2.}
\item
Suppose \(f_n \to f\) in measure. Prove the following:
\begin{examsubparts}
\item[\textnormal{a)}]
If \(f_n \geq 0\) for all~\(n\), then \(\int f \leq \liminf \int f_n\).
\item[\textnormal{b)}]
If \(|f_n| \leq g\) for all~\(n\) and \(g \in L^1\), then \(\int f = \lim \int f_n\).
\end{examsubparts}
\def\ProblemHeadingText{Problem 3.}
\item
Let \((X,\mathcal M,\mu)\) be a~\(\sigma\)-finite measure space, \(\mathcal N\) a sub-\(\sigma\)-algebra of \(\mathcal M\), and \(\nu=\mu|_{\mathcal N}\). Prove that if \(f\in L^1(\mu)\), then there exists \(g\in L^1(\nu)\) (unique modulo~\(\nu\)-null sets) such that
\[
\int_E f\,d\mu=\int_E g\,d\nu
\]
 for all \(E\in\mathcal N\).
\def\ProblemHeadingText{Problem 4.}
\item
Suppose \(1<p<\infty\), \(f\in L^p(0,\infty)\), and \(p^{-1}+q^{-1}=1\). Define
\[
F(x)=\int_0^x f(t)\,dt.
\]
 Show that \(\dfrac{F(x)}{x^{1/q}}\to0\) as \(x\to0\) and \(x\to\infty\).
\def\ProblemHeadingText{Problem 5.}
\item
Let \(\mathcal X,\mathcal Y\) be Banach spaces. Suppose that \((T_n)\) is a sequence of bounded linear operators from \(\mathcal X\) to \(\mathcal Y\) such that \(\lim T_nx\) exists for every \(x\in\mathcal X\). Prove that
\[
Tx:=\lim T_nx
\]
 defines a bounded linear operator from \(\mathcal X\) to \(\mathcal Y\).
\def\ProblemHeadingText{Problem 6.}
\item
Let~\(\mu\) be a~\(\sigma\)-finite positive measure and \(1\leq p\leq\infty\). Let \(g\in L^\infty\). Prove that the operator
\[
Tf=gf
\]
 is bounded on \(L^p\), and \(\lVert T\rVert=\lVert g\rVert_\infty\).
\def\ProblemHeadingText{Problem 7.}
\item
Let \(\mathcal X\) be a normed vector space over \(\mathbb C\). Prove that if \(\mathcal M\) is a closed subspace of \(\mathcal X\) and \(x\in\mathcal X\setminus\mathcal M\), then the subspace \(\mathcal M+\mathbb Cx\) is closed.
\def\ProblemHeadingText{Problem 8.}
\item
Let \(\mathcal E\) be a subset of a Hilbert space \(\mathcal H\). Prove that \((\mathcal E^\perp)^\perp\) is equal to the smallest closed subspace of \(\mathcal H\) containing \(\mathcal E\).
\end{examproblems}
\end{document}
